% TOP_PROBLEM: {"id": 3390, "title": "Wall-Sun-Sun primes and Fibonacci divisibility", "category_id": 1, "category": {"id": 1, "name": "number_theory", "display_name": "Number Theory", "description": "Properties of integers, prime numbers, Diophantine equations.", "slug": "number-theory", "order_index": 1, "created_at": "2026-07-31T15:26:25.670Z"}, "status": "open", "problem_number": "OPG-822", "set_id": 12}
% FIRST_POSTED: 2026-10-01
% ENTRY_KIND: statement_only
% UnsolvedMath ID 3390; source catalogue code OPG-822
% !TeX program = lualatex
% Definition and sources only; this file is not an LLM solution attempt.
\documentclass[11pt,a4paper]{article}
\usepackage[margin=25mm]{geometry}
\usepackage{fontspec}
\setmainfont{lmroman10-regular.otf}[BoldFont=lmroman10-bold.otf,
 ItalicFont=lmroman10-italic.otf,BoldItalicFont=lmroman10-bolditalic.otf]
\usepackage{amsmath,amssymb,mathtools}
\usepackage{unicode-math}
\setmathfont{latinmodern-math.otf}
\AtBeginDocument{\let\setminus\smallsetminus}
\usepackage{xurl,hyperref}
\hypersetup{colorlinks=true,urlcolor=blue}
\setcounter{secnumdepth}{0}
\setlength{\parindent}{0pt}
\setlength{\parskip}{5pt}
\setlength{\emergencystretch}{3em}
\begin{document}
\section*{Wall-Sun-Sun primes and Fibonacci divisibility}\label{3390}
{\small UnsolvedMath ID 3390 · OPG-822 · Number Theory · OpenGarden}
\subsection{Definitions and mathematical statement}
Define the Fibonacci sequence by $F_0=0$, $F_1=1$, and
$F_{m+2}=F_{m+1}+F_m$ for $m\ge0$. For a prime $p$ and nonzero
integer $a$, let $v_p(a)$ be the largest exponent $e\ge0$ such that
$p^e\mid a$. The first formulation of the conjecture is
\[
 \forall\text{ primes }p\ \exists m\ge1\qquad v_p(F_m)=1.
\]
The zero term $F_0$ is excluded. The standard equivalent formulation
for primes $p>5$ is
\[
                  p^2\nmid F_{p-(p/5)}.
\]
Here $(p/5)$ is the Legendre symbol: it is $1$ if the nonzero residue
of $p$ is a square modulo five and $-1$ otherwise. For $p>5$ the
integer $p$ does divide this Fibonacci number, so the nondivisibility
by $p^2$ says its valuation is exactly one. The primes two, three
and five are covered by $F_3=2$, $F_4=3$, and $F_5=5$.

A prime $p>5$ violating the second formula is a Wall--Sun--Sun
prime. The question is their nonexistence, not a finite search-bound
claim. A counterexample must verify the squared-prime divisibility
at the specified index.
\subsection{Short English statement}
For every prime, is some positive Fibonacci number divisible by that prime exactly once?
\subsection{Sources}
\begin{itemize}
\item[S1] ULAM.AI and the UnsolvedMath Contributors, supplied problems.json, record 3390 (OPG-822), OpenGarden collection. Catalogue identity and source transcription; CC BY 4.0.
\url{https://huggingface.co/datasets/ulamai/UnsolvedMath}.
\item[S2] Open Problem Garden, Wall-Sun-Sun primes and Fibonacci divisibility. Original problem and discussion; the mathematical formulation below is expanded from the supplied source record.
\url{http://www.openproblemgarden.org/op/fibonacci_divisibility}.
\end{itemize}
\end{document}
